% !TeX program = lualatex
% Definition and sources only; this file is not an LLM solution attempt.
\documentclass[11pt,a4paper]{article}
\usepackage[margin=25mm]{geometry}
\usepackage{fontspec}
\setmainfont{lmroman10-regular.otf}[BoldFont=lmroman10-bold.otf,
 ItalicFont=lmroman10-italic.otf,BoldItalicFont=lmroman10-bolditalic.otf]
\usepackage{amsmath,amssymb,mathtools}
\usepackage{unicode-math}
\setmathfont{latinmodern-math.otf}
\AtBeginDocument{\let\setminus\smallsetminus}
\usepackage{xurl,hyperref}
\hypersetup{colorlinks=true,urlcolor=blue}
\setcounter{secnumdepth}{0}
\setlength{\parindent}{0pt}
\setlength{\parskip}{5pt}
\setlength{\emergencystretch}{3em}
\begin{document}
\section*{Reflexivity of Point Sets}\label{statement}

\subsection{Definitions and mathematical statement}
Let $S\subset\mathbb R^2$ consist of $n\ge3$ points with no three
collinear. A polygonization is a simple polygon whose vertex set
is exactly $S$. A vertex is reflex when its interior angle is
strictly greater than $\pi$. Write $r(P)$ for the number of
reflex vertices of a polygonization $P$, and define
\[
 \rho(S)=\min_{P:\ V(P)=S}r(P),\qquad
 \rho(n)=\max_{S:\ |S|=n}\rho(S).
\]
The inner minimum chooses the best cyclic order for the particular
point set, while the outer maximum chooses the most difficult
configuration. Both quantifiers are necessary.

Determine $\rho(n)$ exactly, or obtain sharp asymptotic bounds.
The source conjectures the exact expression
$\rho(n)=\lfloor n/4\rfloor$. Establishing it requires both a
polygonization with at most this many reflex corners for every
admissible $S$, and point sets forcing that many corners.

All input points must appear as vertices, including points inside
their convex hull. Replacing $S$ by its hull would remove the
difficulty and is not allowed. The count measures concave corners,
not total angle deficit, and different polygonizations may have
different reflex vertices even on the same set.
\subsection{Short English statement}
What is the largest unavoidable number of concave corners when
a simple polygon must pass through every point of a prescribed
planar set? Is it exactly $\lfloor n/4\rfloor$?
\subsection{Sources}
\begin{itemize}
\item[S1] ULAM.AI and the UnsolvedMath Contributors, supplied problems.json, catalogue entry, AMR Open Problem Lists. Catalogue identity and source transcription; CC BY 4.0.
\url{https://huggingface.co/datasets/ulamai/UnsolvedMath}.
\item[S2] Open Problems Project - Computational Geometry, Problem 66. 
\url{https://topp.openproblem.net/p66}.
\end{itemize}
\end{document}
